\begin{lemma}
\label{editorial-source-lemma-Zariski-topology}
Let $R$ be a ring.
\begin{enumerate}
\item The spectrum of a ring $R$ is empty if and only if $R$
is the zero ring.
\item Every nonzero ring has a maximal ideal.
\item Every nonzero ring has a minimal prime ideal.
\item Given an ideal $I \subset R$ and a prime ideal
$I \subset \mathfrak p$ there exists a prime
$I \subset \mathfrak q \subset \mathfrak p$ such
that $\mathfrak q$ is minimal over $I$.
\item If $T \subset R$, and if $(T)$ is the ideal generated by
$T$ in $R$, then $V((T)) = V(T)$.
\item If $I$ is an ideal and $\sqrt{I}$ is its radical,
see basic notion (\ref{algebra-item-radical-ideal}), then $V(I) = V(\sqrt{I})$.
\item Given an ideal $I$ of $R$ we have $\sqrt{I} =
\bigcap_{I \subset \mathfrak p} \mathfrak p$.
\item If $I$ is an ideal then $V(I) = \emptyset$ if and only
if $I$ is the unit ideal.
\item If $I$, $J$ are ideals of $R$ then $V(I) \cup V(J) =
V(I \cap J)$.
\item If $(I_a)_{a\in A}$ is a set of ideals of $R$ then
$\bigcap_{a\in A} V(I_a) = V(\bigcup_{a\in A} I_a)$.
\item If $f \in R$, then $D(f) \amalg V(f) = \Spec(R)$.
\item If $f \in R$ then $D(f) = \emptyset$ if and only if $f$
is nilpotent.
\item If $f = u f'$ for some unit $u \in R$, then $D(f) = D(f')$.
\item If $I \subset R$ is an ideal, and $\mathfrak p$ is a prime of
$R$ with $\mathfrak p \not\in V(I)$, then there exists an $f \in I$
such that $\mathfrak p \in D(f)$, and $D(f) \cap V(I) = \emptyset$.
\item If $f, g \in R$, then $D(fg) = D(f) \cap D(g)$.
\item If $f_i \in R$ for $i \in I$, then
$\bigcup_{i\in I} D(f_i)$ is the complement of $V(\{f_i \}_{i\in I})$
in $\Spec(R)$.
\item If $f \in R$ and $D(f) = \Spec(R)$, then $f$ is a unit.
\end{enumerate}
\end{lemma}

\begin{proof}
We address each part in the corresponding item below.
\begin{enumerate}
\item This is a direct consequence of (2) or (3).
\item Let $\mathfrak{A}$ be the set of all proper ideals of $R$. This set is
ordered by inclusion and is non-empty, since $(0) \in \mathfrak{A}$ is a proper
ideal. An empty totally ordered subset has the upper bound $(0)$.
Let $A$ be a nonempty totally ordered subset of $\mathfrak A$.
Then $\bigcup_{I \in A} I$ is in
fact an ideal. Since $1 \notin I$ for all $I \in A$, the union does not contain
$1$ and thus is proper. Hence $\bigcup_{I \in A} I$ is in $\mathfrak{A}$ and is
an upper bound for the set $A$. Thus by Zorn's lemma $\mathfrak{A}$ has a
maximal element, which is the sought-after maximal ideal.
\item Since $R$ is nonzero, it contains a maximal ideal which is a prime ideal.
Thus the set $\mathfrak{A}$ of all prime ideals of $R$ is nonempty.
$\mathfrak{A}$ is ordered by reverse-inclusion. An empty totally
ordered subset has any one of these prime ideals as an upper bound.
Let $A$ be a nonempty totally ordered subset of $\mathfrak{A}$.
The intersection $J = \bigcap_{I \in A} I$ is an ideal, and it is proper
because it is contained in any chosen member of $A$.
To prove that it is prime, let $xy \in J$.
Then $xy \in I$ for all $I \in A$. Now let $B = \{I \in A | y \in I\}$. Let $K
= \bigcap_{I \in B} I$. Since $A$ is totally ordered, either $K = J$ (and we're
done, since then $y \in J$) or $K \supset J$ and for all $I \in A$ such that
$I$ is properly contained in $K$, we have $y \notin I$. But that means that for
all those $I$, we have $x \in I$, since they are prime.
Here is why this implies $x \in J$. Since $K\ne J$, choose
$a\in K\setminus J$ and $I_0\in A$ with $a\notin I_0$.
For each $H\in B$, total ordering forces $I_0\subset H$:
the other possibility $H\subset I_0$ would put $a\in K\subset H$
inside $I_0$. Consequently $I_0\subset K$, strictly because
$a\in K\setminus I_0$. This also holds when $B$ is empty,
with its intersection interpreted as $R$.
For any $L\in A$, either $L\subset I_0$, in which case
$L\subsetneq K$ and the preceding primality argument gives $x\in L$,
or $I_0\subset L$, in which case $x\in I_0\subset L$.
Thus $x$ belongs to every $L\in A$ and hence to $J$. In either case,
$J$ is prime as desired. Hence by Zorn's lemma we get a maximal element which
in this case is a minimal prime ideal.
\item This is the same exact argument as (3) except you only consider prime
ideals contained in $\mathfrak{p}$ and containing $I$.
This collection is nonempty because it contains $\mathfrak p$,
which also supplies an upper bound for the empty chain in reverse inclusion.
For a nonempty chain its intersection still contains $I$, is contained
in $\mathfrak p$, and is prime by the argument in (3).
The resulting minimal member $\mathfrak q$ is minimal over $I$
among all primes: any prime between $I$ and $\mathfrak q$ also
lies inside $\mathfrak p$ and therefore belongs to the same collection.
\item $(T)$ is the smallest ideal containing $T$. Hence if $T \subset I$, some
ideal, then $(T) \subset I$ as well. Hence if $I \in V(T)$, then $I \in V((T))$
as well. The other inclusion is obvious.
\item Since $I \subset \sqrt{I}, V(\sqrt{I}) \subset V(I)$. Now let
$\mathfrak{p} \in V(I)$. Let $x \in \sqrt{I}$. Then $x^n \in I$ for some $n$.
Hence $x^n \in \mathfrak{p}$. But since $\mathfrak{p}$ is prime, a boring
induction argument gets you that $x \in \mathfrak{p}$. Hence $\sqrt{I} \subset
\mathfrak{p}$ and $\mathfrak{p} \in V(\sqrt{I})$.
\item Let $f \in R \setminus \sqrt{I}$. Then $f^n \notin I$ for all $n$. Hence
$S = \{1, f, f^2, \ldots\}$ is a multiplicative subset, not containing $0$.
Moreover, $S\cap I=\emptyset$ and $S^{-1}I$ is proper.
Indeed, if $1=a/s$ with $a\in I$ and $s\in S$, equality of fractions
gives $u(s-a)=0$ for some $u\in S$. Then $us=ua\in I\cap S$,
a contradiction. Thus $(S^{-1}R)/(S^{-1}I)$ is a nonzero ring.
By (2) it has a maximal ideal; its inverse image supplies a
prime ideal $\bar{\mathfrak{p}} \subset S^{-1}R$ containing $S^{-1}I$. Then the
pull-back $\mathfrak{p}$ in $R$ of $\bar{\mathfrak{p}}$ is a prime ideal
containing $I$ that does not intersect $S$. This shows that $\bigcap_{I \subset
\mathfrak p} \mathfrak p \subset \sqrt{I}$. Now if $a \in \sqrt{I}$, then $a^n
\in I$ for some $n$. Hence if $I \subset \mathfrak{p}$, then $a^n \in
\mathfrak{p}$. But since $\mathfrak{p}$ is prime, we have $a \in \mathfrak{p}$.
Thus the equality is shown. If $I=R$, both sides are $R$,
with the intersection over the empty set of prime ideals interpreted as $R$.
\item $I$ is not the unit ideal if and only if $I$
is contained in some maximal ideal (to
see this, apply (2) to the ring $R/I$) which is therefore prime.
\item If $\mathfrak{p} \in V(I) \cup V(J)$, then $I \subset \mathfrak{p}$ or $J
\subset \mathfrak{p}$ which means that $I \cap J \subset \mathfrak{p}$. Now if
$I \cap J \subset \mathfrak{p}$, then $IJ \subset \mathfrak{p}$ and hence
either $I \subset \mathfrak{p}$ or $J \subset \mathfrak{p}$, since $\mathfrak{p}$ is prime.
\item $\mathfrak{p} \in \bigcap_{a \in A} V(I_a) \Leftrightarrow
I_a \subset \mathfrak{p}, \forall a \in A \Leftrightarrow
\mathfrak{p} \in V(\bigcup_{a\in A} I_a)$
\item If $\mathfrak{p}$ is a prime ideal and $f \in R$, then either $f \in
\mathfrak{p}$ or $f \notin \mathfrak{p}$ (strictly) which is what the disjoint
union says.
\item If $a \in R$ is nilpotent, then $a^n = 0$ for some $n$. Hence $a^n \in
\mathfrak{p}$ for any prime ideal. Thus $a \in \mathfrak{p}$ as can be shown by
induction and $D(a) = \emptyset$. Now, as shown in (7), if $a \in R$ is not
nilpotent, then there is a prime ideal that does not contain it.
\item $f \in \mathfrak{p} \Leftrightarrow uf \in \mathfrak{p}$, since $u$ is
invertible.
\item If $\mathfrak{p} \notin V(I)$, then $\exists f \in I \setminus
\mathfrak{p}$. Then $f \notin \mathfrak{p}$ so $\mathfrak{p} \in D(f)$. Also if
$\mathfrak{q} \in D(f)$, then $f \notin \mathfrak{q}$ and thus $I$ is not
contained in $\mathfrak{q}$. Thus $D(f) \cap V(I) = \emptyset$.
\item If $fg \in \mathfrak{p}$, then $f \in \mathfrak{p}$ or $g \in
\mathfrak{p}$. Hence if $f \notin \mathfrak{p}$ and $g \notin \mathfrak{p}$,
then $fg \notin \mathfrak{p}$. Since $\mathfrak{p}$ is an ideal, if $fg \notin
\mathfrak{p}$, then $f \notin \mathfrak{p}$ and $g \notin \mathfrak{p}$.
\item $\mathfrak{p} \in \bigcup_{i \in I} D(f_i) \Leftrightarrow \exists i \in
I, f_i \notin \mathfrak{p} \Leftrightarrow \mathfrak{p} \in \Spec(R)
\setminus V(\{f_i\}_{i \in I})$
\item If $D(f) = \Spec(R)$, then $V(f) = \emptyset$ and
hence $fR = R$, so $f$ is a unit.
\end{enumerate}
\end{proof}